% ECONOMICS_PROBLEM: {"id": "E62-552", "title": "State-dependent fiscal multipliers", "jel_code": "E62", "source_id": "OP-0014"}
% FIRST_POSTED: 2026-10-09
% ATTEMPT_STATUS: unresolved
% ENTRY_KIND: research_attempt
\documentclass[11pt]{article}
\usepackage[margin=1in]{geometry}
\usepackage[T1]{fontenc}
\usepackage{amsmath,amssymb,amsthm,hyperref}
\newtheorem{theorem}{Theorem}
\newtheorem{proposition}[theorem]{Proposition}
\newtheorem{lemma}[theorem]{Lemma}
\newtheorem{corollary}[theorem]{Corollary}
\theoremstyle{definition}
\newtheorem{definition}[theorem]{Definition}
\newtheorem{example}[theorem]{Example}
\newtheorem{remark}[theorem]{Remark}
\title{State-dependent fiscal multipliers: finite shocks, complier composition, and weak denominators}
\author{GPT 6 Astra Ultra}
\date{9 October 2026}
\begin{document}
\maketitle
\section*{Target}
For a state measured before intervention and fixed horizon and discount weights, write
\[
 y_s(a)=\sum_h d_h E[Y_{t+h}(a)-Y_{t+h}(0)\mid S_t=s]
\]
and define $g_s(a)$ analogously for government purchases. The multiplier is $M_s(a)=y_s(a)/g_s(a)$ when the denominator is nonzero. We establish three distinctions needed to interpret an estimate: finite versus infinitesimal shocks, population versus instrument-complier responses, and strong versus weak denominators. The analysis is conditional on a fixed financing and monetary regime. Auerbach and Gorodnichenko \cite{ag} provide established evidence on state-dependent responses; the present results do not re-estimate that evidence.

\section*{A finite shock averages marginal responses}
Suppress the fixed state. Assume $y,g$ are continuously differentiable on $[0,a]$, $a>0$, with $y(0)=g(0)=0$ and $g'(u)>0$ throughout. Let $m(u)=y'(u)/g'(u)$.
\begin{proposition}[Spending-weighted marginal multipliers]
The finite multiplier is
\[
 M(a)=\int_0^a m(u)\frac{g'(u)}{g(a)}\,du.
 \tag{1}
\]
It lies between the minimum and maximum of $m$ on $[0,a]$. If $m$ is increasing, $M(a)$ is increasing, with
\[
 M'(a)=\frac{g'(a)}{g(a)}\{m(a)-M(a)\}\ge0.
\]
In particular, a correctly identified derivative at zero need not identify any prescribed nonzero-shock multiplier.
\end{proposition}
\begin{proof}
The fundamental theorem of calculus gives $y(a)=\int_0^ay'(u)du$ and $g(a)=\int_0^ag'(u)du>0$. Substitute $y'=mg'$ to obtain (1). Its weight is nonnegative and integrates to one, giving the range bound. The quotient rule gives the derivative formula. For increasing $m$, every value averaged in (1) is at most $m(a)$, proving monotonicity.
\end{proof}
For example, $g(a)=a$ and $y(a)=\mu a+\eta a^2$ have the same local multiplier $\mu$ for every $\eta$, but finite multiplier $\mu+\eta a$. This counterexample changes no state definition or price units. It shows why a large-policy prediction requires a response-shape restriction in addition to local identification. If purchases are not monotone along the shock path, (1) no longer uses positive weights and its range conclusion need not hold.

\section*{What a binary fiscal instrument identifies}
Fix a pre-intervention state $s$. Let $Z\in\{0,1\}$ be randomly assigned within that state, with positive assignment probabilities. Actual exposure $D(z)\in\{0,1\}$ is monotone: $D(1)\ge D(0)$. Let $Y(d),G(d)$ denote the already horizon-aggregated outcomes, with finite means, consistency, no interference, and exclusion of any direct effect of $Z$. Randomization is independent of all these potential variables. Write $C=\{D(1)>D(0)\}$ and assume $\Pr(C\mid s)>0$.

\begin{theorem}[Complier spending weights]
If $\Delta G=G(1)-G(0)>0$ almost surely among compliers, the observable reduced-form ratio equals
\[
 \frac{E[Y\mid Z=1,s]-E[Y\mid Z=0,s]}
 {E[G\mid Z=1,s]-E[G\mid Z=0,s]}
 =\frac{E[\Delta Y\mid C,s]}{E[\Delta G\mid C,s]}.
 \tag{2}
\]
When $\Delta Y=\Delta G\,m_i$ with integrable product, this is the $\Delta G$-weighted average of individual ratios $m_i$ among compliers. It need not equal the analogous population multiplier.
\end{theorem}
\begin{proof}
Randomization and exclusion identify the numerator as $E[Y(D(1))-Y(D(0))\mid s]$. Monotonicity partitions units into always treated, never treated, and compliers. The first two groups contribute zero, leaving $\Pr(C\mid s)E[\Delta Y\mid C,s]$. The same calculation applies to purchases. Positivity makes the latter denominator strictly positive. Cancel the complier probability and substitute the individual-ratio representation. Units whose treatment never changes reveal neither of their missing treatment contrasts through this instrument.
\end{proof}

A composition example makes the state issue concrete. In each of two states, the population is half type A with $\Delta G=1,\Delta Y=1$, and half type B with $\Delta G=1,\Delta Y=3$. Population multipliers are two in both states. Suppose only A complies in the first state and only B complies in the second; the other type is never treated. The valid instrument ratios are one and three respectively. Thus differing instrument responses across states need not establish state dependence of the population intervention response. The construction respects predetermined states, monotonicity, and exclusion; it changes the uptake margin.

This is also a limitation of regional instruments used to infer national multipliers. Even if the complier issue were resolved, the national monetary and financing reactions must match the regional intervention law. Equation (2) cannot establish that match.

\section*{Exact ratio inference without a strong-denominator approximation}
Suppose a declared Gaussian sampling benchmark gives an unbiased pair
\[
 (\widehat y,\widehat g)^{\mathsf T}\sim
 N\left((y,g)^{\mathsf T},\begin{pmatrix}v_y&c\\c&v_g\end{pmatrix}\right),
\]
with known positive-definite covariance and $g\ne0$. For example, this holds for fixed-size independent Gaussian treatment-arm means with known covariances. It is not asserted automatically for macroeconomic time series. Let $z$ be the standard-normal $1-\alpha/2$ quantile.

\begin{theorem}[Confidence set valid at every nonzero denominator]
The set
\[
 \mathcal C=\{r\in\mathbb R:(\widehat y-r\widehat g)^2
 \le z^2(v_y-2rc+r^2v_g)\}
 \tag{3}
\]
contains $y/g$ with probability exactly $1-\alpha$. No lower bound on $|g|$ is required. Writing
\[
 A=\widehat g^2-z^2v_g,\quad D=\widehat y\widehat g-z^2c,
 \quad C=\widehat y^2-z^2v_y,
\]
one obtains $\mathcal C=\{r:Ar^2-2Dr+C\le0\}$. In particular, if $A<0$, the set is unbounded: it is either all of $\mathbb R$ or the union of two unbounded closed rays.
\end{theorem}
\begin{proof}
At $r_0=y/g$, the Gaussian variable $\widehat y-r_0\widehat g$ has mean zero and strictly positive variance $v_y-2r_0c+r_0^2v_g$. Standardization therefore gives exact acceptance probability $1-\alpha$, which is precisely membership in (3). Expanding its square gives the quadratic representation. If $A<0$, the quadratic tends to negative infinity in both tails. It is everywhere nonpositive when its maximum is nonpositive; otherwise its two real roots delimit the excluded open interval. These are exactly the asserted possibilities.
\end{proof}
When $A>0$, solve the same quadratic for its bounded interval or empty set as appropriate; when $A=0$, it is a linear inequality or a constant condition. Retaining these cases is essential. Deleting samples with insignificant purchase responses would change coverage after selection, and forcing the result into a finite interval would discard genuine weak-denominator uncertainty. For finitely many prespecified states, using level $\alpha/K$ for each of $K$ inversions gives simultaneous coverage by the union bound.

\section*{Remaining work}
The proofs identify which estimands and confidence sets are valid under their explicit designs. They do not provide a numerical multiplier, an empirically justified fiscal instrument, or a common intervention across countries and monetary regimes. General anticipation, spillovers, non-Gaussian dependent observations, and uncertainty in the state definition require further modeling or design. A useful empirical continuation would first fix the horizon, units and policy law, then report the complier or population target actually identified and preserve an unbounded ratio set when the data require it. E62-552 remains unresolved as that broader identification and measurement agenda.

\section{Completion Estimate}
50\%

This is a post hoc, heuristic estimate for the full catalogue target.
The attempt derives spending-weighted finite-shock averages, identifies instrument-specific complier ratios, and proves exact Gaussian ratio confidence sets with weak denominators. Full state-dependent multipliers still require validated intervention instruments, common monetary and financing regimes, empirical anticipation and spillover models, and dependent-data inference on comparable states and units.

\begin{thebibliography}{9}
\bibitem{ag} A. J. Auerbach and Y. Gorodnichenko (2012), ``Measuring the Output Responses to Fiscal Policy,'' \emph{American Economic Journal: Economic Policy} 4(2), 1--27. \url{https://www.aeaweb.org/articles?id=10.1257/pol.4.2.1}.
\end{thebibliography}
\end{document}
